\subsection{皮肤可见疾病与性自信：银屑病、湿疹、白癜风与暴露焦虑}

有一类皮肤病的特点不是"长在私处"，而是\textbf{长在别人看得见的地方}——头皮、面部、颈部、手臂、小腿、指甲。银屑病（牛皮癣）在我国的患病率约 0.5\%，湿疹与特应性皮炎影响的人群更广，白癜风患病率约 0.5\%--2\%。它们绝大多数不影响性功能，也不累及生殖器；真正受累的，往往是患者的\textbf{自我展示信心}——而被注视的恐惧，恰恰会绕过生理，直接削弱亲密意愿。

\subsubsection{先说清三件事：不传染、不是性病、可以控制}

\begin{itemize}
\item \textbf{不传染}：银屑病、湿疹、白癜风都不是感染性疾病，任何形式的亲密接触——牵手、拥抱、共浴、性行为——都不会把皮损"传"给伴侣。这是最需要说破、也最常被患者自己默默怀疑的一点。
\item \textbf{不是性病}：这些疾病与性行为没有任何因果关联。但皮损出现在外阴、腹股沟等部位时（银屑病可累及生殖器，称为反向银屑病/屈曲部位银屑病），患者极易陷入"会不会被当成性病"的恐慌——这是就医澄清可解决的问题，详见本节末尾。
\item \textbf{可以控制}：现代皮肤科对三者的控制手段都远比公众印象成熟——外用制剂、光疗、系统药物乃至生物制剂，多数患者可实现皮损大幅清除。"不治之症"的旧印象，是焦虑的重要来源，也是不准确的。
\end{itemize}

\subsubsection{暴露焦虑如何影响亲密关系}

\begin{itemize}
\item 典型的循环是：皮损发作$\rightarrow$怕被看见$\rightarrow$拒绝共浴、关灯、穿遮盖衣物$\rightarrow$回避亲密$\rightarrow$伴侣感到被推远$\rightarrow$关系疏远$\rightarrow$自信进一步下降。打断这个循环的关键动作只有一个：\textbf{把病说出来}。
\item 焦虑常在"新关系初期"最重——脱衣的瞬间被体验为"暴露审判"。而研究身体意象的学者反复确认的一件事是：伴侣对皮损的在意程度，几乎总是远低于患者的想象；他们注意到的是你的紧张，而不是你的皮肤。
\item 长期遮盖与回避还可能被伴侣误读为"性冷淡""不喜欢我"——对方在不知道原因的情况下，只能往自己身上找解释。沉默造成的误会，常常比皮损本身更伤关系。
\end{itemize}

\subsubsection{性生活实务：皮损、药物与体位}

\begin{itemize}
\item \textbf{护理时机}：外用保湿剂或药物后稍等片刻、吸收完毕再开始亲密活动，可减少黏腻感与摩擦刺激；皮损急性期有裂口、渗出时，先处理、再亲密。
\item \textbf{摩擦与受压}：肘、膝、肩胛等好发部位的摩擦可能加重皮损（同形现象主要见于银屑病急性期）。可以换用压力更分散的体位、铺垫软枕，或在易摩擦部位提前使用保湿霜——这些都不影响体验，只是多一分照护。
\item \textbf{药物注意事项}：外用糖皮质激素长期用于腹股沟等薄嫩部位需遵医嘱；部分维 A 酸类药物可能有黏膜干燥；光疗安排与性生活时间错开即可，无需禁忌。系统性用药（甲氨蝶呤、生物制剂等）一般不影响性功能，具体顾虑直接问皮肤科医生。
\item \textbf{把皮肤管理纳入生活节律}：发作期的情绪低落是真实的（银屑病与抑郁的共病率显著升高），此时可借鉴慢性病亲密管理的思路——非插入式亲密、低耗能方案、症状缓解窗口期安排（参见慢性病患者的性健康管理章节）。
\end{itemize}

\subsubsection{如何向伴侣说明：时机与话术}

\begin{itemize}
\item \textbf{时机}：不必在第一次约会就交代病历，但在关系走向亲密之前说清，几乎总是更好的选择——把"被发现"变成"我告诉你"，主动权完全不同。
\item \textbf{话术示例}："我的皮肤有一种慢性问题，叫银屑病，不传染，也不是性病，会反复但可以控制。你想了解的话我可以多说说；亲密的时候如果某个部位你介意，我们调整就好。"
\item \textbf{允许对方有反应时间}：多数伴侣的反应是接受或 indifferent（不在意）；若对方有明显顾虑，可提供科普资料（医院患教材料、皮肤科医生科普），而不必当场辩解。
\end{itemize}

\subsubsection{伴侣可以做什么}

\begin{itemize}
\item 用平常心对待皮损：不盯着看、也不刻意回避看；对方提起时倾听，而不是急于给"根治偏方"建议。
\item 了解基本事实：不传染、慢性、可控制——这三条知道，就能避开绝大多数踩雷发言。
\item 在对方低落时提供具体的帮助：陪伴复诊、帮忙涂背部的药、发作期主动承担更多家务。照顾行动比"别想太多"有用一百倍。
\end{itemize}

\subsubsection{何时就医：容易被误当性病的皮肤问题}

\begin{itemize}
\item 外阴、腹股沟、肛周出现红斑、脱屑、白斑时，建议直接挂皮肤性病科——"性病科"的名字让许多人却步，但它恰恰是鉴别"是不是性病"最专业的科室。一次专业澄清，能换来长年的心安。
\item 需要鉴别的常见情况包括：反向银屑病与股癣、外阴湿疹与接触性皮炎、白癜风与外阴白色病变（女性卷另有专节讨论外阴白色斑块的鉴别）、以及真正的性传播疾病。自行对照网络图片下结论，是焦虑放大器。
\item 若皮损反复出现在生殖器且伴随性交不适，或皮肤问题已明显影响情绪（持续低落、回避社交），请同时寻求皮肤科与心理支持，两条线并行。
\end{itemize}

\noindent\textit{【编者注】}皮肤可见疾病与性自信的议题，核心是一个"暴露"的心理动作，而不是一个皮肤科技术问题。医学信息请以皮肤科面诊为准；本节希望传递的是：皮肤不是爱的入场券，把病说出口的那一刻，暴露焦虑的循环就开始松动了。身体意象的总体讨论见本章前文；残障与身体改变者的亲密议题另有专节。
